% gz-p10-07-2: 切线放缩——y=ln(1+x) 在 x=0 处切线 y=x（红虚），显示 ln(1+x) ≤ x 取等当 x=0
\documentclass[tikz,border=4pt]{standalone}
\usepackage{ctex}
\usepackage{amsmath}
\usepackage{pgfplots}
\pgfplotsset{compat=1.18}

\begin{document}
\begin{tikzpicture}
\begin{axis}[
  width=10cm, height=7cm,
  axis lines=center,
  xlabel={$x$}, ylabel={$y$},
  every axis x label/.style={at={(axis cs:3.6,0)}, anchor=west},
  every axis y label/.style={at={(axis cs:0,3.7)}, anchor=south},
  xmin=-1.0, xmax=3.6,
  ymin=-2.0, ymax=3.7,
  xtick={-0.5, 1, 2, 3},
  ytick={-0.5, 1, 2, 3},
  tick style={color=black},
  ticklabel style={font=\small},
  clip=true,
]
  % 曲线 y = ln(1+x)
  \addplot[blue, thick, domain=-0.95:3.5, samples=120] {ln(1+x)};
  \node[right, blue] at (axis cs:2.4, 1.0) {$y=\ln(1+x)$};

  % 切线 y = x（红虚），过原点
  \addplot[red, dashed, thick, domain=-1.0:3.5] {x};
  \node[right, red] at (axis cs:1.7, 2.1) {切线 $y=x$};

  % 切点 (0, 0)
  \addplot[only marks, mark=*, mark size=3pt, red] coordinates {(0,0)};
  \node[above left, red, font=\small] at (axis cs:0,0) {切点};

  % 注释：ln(1+x) ≤ x，取等当 x = 0
  \node[font=\footnotesize, align=left] at (axis cs:1.5, -1.4)
    {$\ln(1+x)\le x$\\等号成立 $\Longleftrightarrow x=0$};

  % 强调蓝总在红下方
  \draw[gray, dashed] (axis cs:2.0, 2.0) -- (axis cs:2.0, 1.0986);
  \node[right, font=\footnotesize, gray] at (axis cs:2.0, 1.55) {差};
\end{axis}
\end{tikzpicture}
\end{document}
